% ===========================================================================
%  preamble.tex  --  packages and global settings (compile with XeLaTeX)
%  Normally you do not need to edit this file; thesis data goes in main.tex.
% ===========================================================================

% --- Math and fonts ---------------------------------------------------------
\usepackage{amsmath, amssymb}
\usepackage{unicode-math}                 % loads fontspec
\defaultfontfeatures{Ligatures=TeX}
% Fonts are loaded by FILE name (found through kpathsea), so this works on any
% TeX Live / MiKTeX / Overleaf, even when the fonts are not installed in the OS.
% Main text: Adobe Caslon Pro (files in fonts/). Old-style figures in running text.
\IfFileExists{fonts/ACaslonPro-Regular.otf}{%
\setmainfont{ACaslonPro}[
  Path           = fonts/ ,
  Extension      = .otf ,
  UprightFont    = *-Regular ,
  BoldFont       = *-Bold ,
  ItalicFont     = *-Italic ,
  BoldItalicFont = *-BoldItalic ,
  Ligatures      = {TeX, Common} ,
  Numbers        = OldStyle ,
  SmallCapsFeatures = {Letters=SmallCaps} ]
}{%  FALLBACK (free font) when the Caslon files are not in fonts/
\setmainfont{texgyrepagella-regular.otf}[
  BoldFont = texgyrepagella-bold.otf , ItalicFont = texgyrepagella-italic.otf ,
  BoldItalicFont = texgyrepagella-bolditalic.otf ,
  Ligatures = {TeX, Common} , Numbers = OldStyle ,
  SmallCapsFeatures = {Letters=SmallCaps} ]}
% Headings: Caslon Semibold (old-style figures, same as the running text)
\IfFileExists{fonts/ACaslonPro-Regular.otf}{%
\newfontfamily\headingfont{ACaslonPro}[
  Path           = fonts/ ,
  Extension      = .otf ,
  UprightFont    = *-Semibold ,
  BoldFont       = *-Bold ,
  ItalicFont     = *-SemiboldItalic ,
  BoldItalicFont = *-BoldItalic ,
  Ligatures      = {TeX, Common} ,
  Numbers        = OldStyle ]
}{%
\newfontfamily\headingfont{texgyrepagella-bold.otf}[
  BoldFont = texgyrepagella-bold.otf , ItalicFont = texgyrepagella-bolditalic.otf ,
  Ligatures = {TeX, Common} , Numbers = OldStyle ]}
% Sans and monospace (used rarely) and math stay Latin Modern (loaded by file name)
\setsansfont{lmsans10-regular.otf}[
  BoldFont       = lmsans10-bold.otf,
  ItalicFont     = lmsans10-oblique.otf,
  BoldItalicFont = lmsans10-boldoblique.otf]
\setmonofont{lmmono10-regular.otf}[
  ItalicFont     = lmmono10-italic.otf,
  BoldFont       = lmmonolt10-bold.otf,
  BoldItalicFont = lmmonolt10-boldoblique.otf]
\setmathfont{latinmodern-math.otf}
% Alternative look (needs the TeX Gyre fonts, included in a full TeX Live):
%   \setmainfont{TeX Gyre Pagella}   \setsansfont{TeX Gyre Heros}[Scale=MatchLowercase]
%   \setmonofont{TeX Gyre Cursor}[Scale=MatchLowercase]  \setmathfont{TeX Gyre Pagella Math}
% Also: Libertinus Serif / Libertinus Math, STIX Two Text / STIX Two Math.
% Digits in math mode taken from Caslon (lining figures), so $-0.6$ or $G=21$ match the
% text and siunitx numbers. Letters and symbols stay Latin Modern: Caslon is a text font
% without math metrics (Caslon italic letters in math give poor spacing, e.g. f(x), K_s).
\IfFileExists{fonts/ACaslonPro-Regular.otf}{%
  \setmathfont{ACaslonPro-Regular.otf}[Path=fonts/, Numbers=Lining, range={up/num}]}{}

% --- Languages: English (main) + Spanish -----------------------------------------
\usepackage{polyglossia}
\setmainlanguage[variant=american]{english}
\setotherlanguage{spanish}
\usepackage{csquotes}                     % \enquote{...} adapts to the active language

% --- Page layout and spacing -------------------------------------------------------
\usepackage[left=3cm, right=2.5cm, top=2.5cm, bottom=2.5cm, headheight=18pt]{geometry}
\usepackage{setspace}
\onehalfspacing                           % 1.5 line spacing (floats/footnotes stay single)
\usepackage{microtype}
\setlength{\emergencystretch}{2em}        % avoids slightly overfull lines (long file names, code)
\usepackage{enumitem}
\setlist{nosep}

% --- Colours ---------------------------------------------------------------------------
\usepackage[dvipsnames,table]{xcolor}
\definecolor{uablue}{HTML}{0B3C8C}
\definecolor{mdtRed}{rgb}{0.5,0,0}

% --- Headings, headers and footers (KOMA-Script) -------------------------------
\KOMAoptions{chapterprefix=true, listof=totoc, bibliography=totoc}
\setkomafont{disposition}{\headingfont}
\setkomafont{chapter}{\headingfont\Huge\color{uablue}}
\setkomafont{chapterprefix}{\headingfont\large\color{uablue}}
\setkomafont{section}{\headingfont\Large\color{uablue}}
\setkomafont{subsection}{\headingfont\large}
\usepackage[automark]{scrlayer-scrpage}
\clearpairofpagestyles
\ihead{\headmark}                         % left:  chapter title
\ohead{\pagemark}                         % right: page number
\cfoot[\pagemark]{}                       % first page of a chapter: centred number
\setkomafont{pageheadfoot}{\small\itshape}
\setkomafont{pagenumber}{\small\upshape}
\KOMAoptions{headsepline=true}
\addtokomafont{caption}{\small}
\addtokomafont{captionlabel}{\bfseries}

% --- Figures --------------------------------------------------------------------------------
\usepackage{graphicx}
\usepackage{caption}
\usepackage{subcaption}
\captionsetup{font=small, labelfont=bf}
\captionsetup[table]{position=above}
\captionsetup[longtable]{position=above, labelfont=normalfont}   % same look as tabularray captions      % table captions above (astronomy convention)
\captionsetup[figure]{position=below}
\usepackage{rotating}                     % sidewaysfigure / sidewaystable

% Less strict float placement (helps on chapter-opening pages)
\renewcommand{\topfraction}{0.85}
\renewcommand{\bottomfraction}{0.7}
\renewcommand{\textfraction}{0.15}
\renewcommand{\floatpagefraction}{0.7}

% One figure folder per chapter: figures/chapter1/, figures/chapter2/, ...
% Call \chapterfigs{chapterN} at the top of each chapter file.
\newcommand{\chapterfigs}[1]{\graphicspath{{figures/#1/}}}

% --- Tables: tabularray + booktabs-style rules + siunitx --------------------------
\usepackage{tabularray}
\UseTblrLibrary{booktabs, siunitx}
\usepackage{booktabs}
\usepackage{multirow}
\usepackage{pdflscape}                    % landscape pages for large tables
\usepackage{longtable}                    % fast multi-page tables (see \uahead and longcompact below)

\definecolor{tablegray}{HTML}{F2F4F7}     % light fill: zebra rows / highlighted column
\definecolor{okgreen}{HTML}{1E8449}
\definecolor{warnorange}{HTML}{B9770E}
\definecolor{badred}{HTML}{B03A2E}

% Lining figures for tables (running text keeps the old-style figures of Caslon)
% Declared ONCE as a font family (cheap to select). Do not use \addfontfeatures inside
% table cells: tabularray typesets every cell several times and compilation becomes very slow.
\IfFileExists{fonts/ACaslonPro-Regular.otf}{%
\newfontfamily\tablenums{ACaslonPro}[
  Path           = fonts/ ,
  Extension      = .otf ,
  UprightFont    = *-Regular ,
  BoldFont       = *-Bold ,
  ItalicFont     = *-Italic ,
  BoldItalicFont = *-BoldItalic ,
  Ligatures      = {TeX, Common} ,
  Numbers        = Lining ,
  SmallCapsFeatures = {Letters=SmallCaps} ]
}{%
\newfontfamily\tablenums{texgyrepagella-regular.otf}[
  BoldFont = texgyrepagella-bold.otf , ItalicFont = texgyrepagella-italic.otf ,
  BoldItalicFont = texgyrepagella-bolditalic.otf ,
  Ligatures = {TeX, Common} , Numbers = Lining ,
  SmallCapsFeatures = {Letters=SmallCaps} ]}
\newcommand{\tabfont}{\small\tablenums}
\newcommand{\tabfontsmall}{\footnotesize\tablenums}   % compact tables

% Thesis table style: plain white header with bold text, thin booktabs-like rules.
%   uatblr      plain table                        (like tblr)
%   uatalltblr  table with notes, inside \table     (like talltblr)
%   ualongtblr  multi-page table, header repeated   (like longtblr)
\NewTblrEnviron{uatblr}
\NewTblrEnviron{uatalltblr}
\NewTblrEnviron{ualongtblr}
\SetTblrOuter[uatalltblr]{tall}
\SetTblrOuter[ualongtblr]{long}
\SetTblrInner[uatblr,uatalltblr,ualongtblr]{
  rowsep   = 2pt,
  colsep   = 6pt,
  cells    = {font=\tabfont},
  row{1}   = {font=\tabfont\bfseries},
  hline{1,Z} = {0.09em},
  hline{2}   = {0.05em},
}
% Compact profile for long tables (more rows per page): smaller type, tighter rows
\SetTblrInner[ualongtblr]{
  rowsep   = 1pt,
  colsep   = 4pt,
  cells    = {font=\tabfontsmall},
  row{1}   = {font=\tabfontsmall\bfseries},
}
% Tables are typeset with single line spacing (the thesis body is 1.5-spaced)
\usepackage{etoolbox}
\AtBeginEnvironment{ualongtblr}{\singlespacing}
\SetTblrStyle{remark}{font=\footnotesize\tablenums}
\SetTblrStyle{note}{font=\footnotesize\tablenums}

% FAST long tables: plain longtable + booktabs (much faster than longtblr, same look).
%   \begin{longcompact} ... \end{longcompact}  wraps a longtable: small type, lining
%   figures, single spacing, plain white header with thin rules.   \uahead{text} = bold header text.
%   For zebra rows add \rowcolors{2}{tablegray}{white} inside longcompact.
\newcommand{\uahead}[1]{\textbf{#1}}
\newenvironment{longcompact}
  {\begingroup\footnotesize\tablenums\singlespacing
   \setlength{\tabcolsep}{6pt}\renewcommand{\arraystretch}{1.15}}
  {\endgroup}

% Status words (coloured, bold) for use inside tables
\newcommand{\statusok}[1]{\textcolor{okgreen}{\textbf{#1}}}
\newcommand{\statuswarn}[1]{\textcolor{warnorange}{\textbf{#1}}}
\newcommand{\statusbad}[1]{\textcolor{badred}{\textbf{#1}}}

% --- Units and numbers: \num{1.2e4}, \qty{3}{\kilo\parsec}, \ang{;;1.5} ------------
\usepackage{siunitx}
\sisetup{
  mode                 = text,       % numbers and units in the text font (Caslon), not in math mode
  reset-text-family    = false,      % inherit the current font (bold, italic...)
  text-font-command    = {\tablenums}, % lining (equal-height) figures everywhere, also in running text
  group-digits         = integer,
  range-phrase         = {--},
  range-units          = single,
  per-mode             = symbol,
}
\DeclareSIUnit\parsec{pc}
\DeclareSIUnit\solarmass{\text{M}\ensuremath{_{\odot}}}
\DeclareSIUnit\solarlum{\text{L}\ensuremath{_{\odot}}}
\DeclareSIUnit\year{yr}
\DeclareSIUnit\Gyr{Gyr}
\DeclareSIUnit\Myr{Myr}

% --- Astronomy: journal macros (\apj, \mnras, \aap ...) used by ADS BibTeX ---
\usepackage{aas_macros}
\makeatletter\let\jnl@style\upshape\makeatother   % aas_macros uses \rm, which KOMA-Script does not define

% --- Bibliography: natbib + A&A/astronomy style (aa.bst) -------------------------
\usepackage[authoryear, round]{natbib}
\bibpunct{(}{)}{;}{a}{}{,}                % (Kennicutt 1998; Speagle et al. 2014)
\setlength{\bibsep}{0.35em}     % space between entries (default is much larger at 1.5 spacing)
\setlength{\bibhang}{1.5em}    % hanging indent for entries
% Lining figures for years in citations and for the whole reference list
\makeatletter
\apptocmd{\NAT@parse}{\begingroup\edef\x{\endgroup\noexpand\gdef\noexpand\NAT@date{{\noexpand\tablenums\NAT@date}}}\x}{}{\PackageWarning{preamble}{natbib patch failed}}
\makeatother
% Appendices: TOC entries read "Appendix A  Title" (automatic, nothing to add in the appendix files)
\newcommand*{\uaappendixtoc}{%
  \renewcommand*{\addchaptertocentry}[2]{\ifstrempty{##1}{\addtocentrydefault{chapter}{}{##2}}{\addtocentrydefault{chapter}{}{\appendixname\nobreakspace ##1\texorpdfstring{\enskip}{ }##2}}}}
\apptocmd{\appendix}{\uaappendixtoc}{}{\PackageWarning{preamble}{appendix patch failed}}
% Reference list: slightly smaller type, single spacing (the body is 1.5-spaced)
\AtBeginEnvironment{thebibliography}{\small\tablenums\singlespacing}

% --- Hyperlinks (load after natbib; keep cleveref after hyperref) --------------
\usepackage[unicode, pdfusetitle=false]{hyperref}
\hypersetup{
  colorlinks   = true,                    % set to false for the printed version
  linkcolor    = uablue,
  citecolor    = uablue,
  urlcolor     = uablue,
  bookmarksopen= true,
  bookmarksopenlevel = 1,
  pdfstartview = FitV,
  breaklinks   = true,
}
\usepackage[capitalise, nameinlink]{cleveref}   % \cref{fig:x} -> "Figure 1"

% --- Acronyms (glossaries-extra; definitions in acronyms.tex) --------------------
%   The list is built automatically (option automake + rule in latexmkrc).
%   automake       : builds the list itself during compilation (works on Overleaf)
%   nomain         : no general glossary, only the list of acronyms
%   toc            : list appears in the table of contents
%   nonumberlist   : no page numbers in the list (remove to list the pages where used)
%   nohypertypes   : acronyms in the text are NOT coloured links (delete for links)
\usepackage[automake, acronym, nomain, toc, nonumberlist, nopostdot, nogroupskip,
            nohypertypes={acronym}]{glossaries-extra}
\setabbreviationstyle[acronym]{long-short}      % first use: "Milky Way (MW)"
\makeglossaries
\setglossarystyle{long}                         % two columns: acronym | meaning
\setlength{\glsdescwidth}{0.78\textwidth}
\renewcommand*{\glsnamefont}[1]{{\tablenums\textbf{#1}}}   % bold, lining figures (2MASS, PS1)
% Astronomy style first use with a citation:
%   \glscite{2MASS}{Skrutskie2006} -> "Two Micron All-Sky Survey (2MASS; Skrutskie et al. 2006)"
%   (afterwards it prints just "2MASS"). Optional 3rd argument in [] = post-note of \citep.
\NewDocumentCommand{\glscite}{m m O{}}{%
  \ifglsused{#1}{\gls{#1}}{%
    \glsxtrlong{#1} \citep[\glsxtrshort{#1};][#3]{#2}\glsunset{#1}}}
% \printacronymlist: prints the "List of Acronyms" (called in main.tex after \listoftables)
\newcommand{\printacronymlist}{%
  \printglossary[type=\acronymtype, title={List of Acronyms}, toctitle={List of Acronyms}]}

% --- Personal macros -------------------------------------------------------------------------
% ---------------------------------------------------------------------------
% Personal macros (astronomy shortcuts). Add your own here.
% ---------------------------------------------------------------------------
\usepackage{xspace}

% Spectroscopic notation: \hi -> H I,  \hii -> H II,  \nii -> [N II]
\newcommand{\hi}{H\,\textsc{i}\xspace}
\newcommand{\hii}{H\,\textsc{ii}\xspace}
\newcommand{\nii}{[N\,\textsc{ii}]\xspace}
\newcommand{\oiii}{[O\,\textsc{iii}]\xspace}
\newcommand{\Halpha}{H$\alpha$\xspace}
\newcommand{\Hbeta}{H$\beta$\xspace}

% Abundances and solar units (math mode)
\newcommand{\FeH}{\ensuremath{[\mathrm{Fe}/\mathrm{H}]}\xspace}
\newcommand{\aFe}{\ensuremath{[\alpha/\mathrm{Fe}]}\xspace}
\newcommand{\Msun}{\ensuremath{\mathrm{M}_{\odot}}\xspace}
\newcommand{\Lsun}{\ensuremath{\mathrm{L}_{\odot}}\xspace}
\newcommand{\Mstar}{\ensuremath{M_{\star}}\xspace}
\newcommand{\SFR}{\ensuremath{\mathrm{SFR}}\xspace}
\newcommand{\sSFR}{\ensuremath{\mathrm{sSFR}}\xspace}
\newcommand{\Sgas}{\ensuremath{\Sigma_{\mathrm{gas}}}\xspace}
\newcommand{\Ssfr}{\ensuremath{\Sigma_{\mathrm{SFR}}}\xspace}

% Reminder note inside the text (shows in red; remove before final version)
\newcommand{\todo}[1]{\textcolor{red}{\textbf{[TODO: #1]}}}
